\section{Evaluation Methodology}
\label{sec:evaluation_methodology}

This section presents the framework used to evaluate the predictive performance of all club-level models on real \textit{Brasileirão Série A} data from 2020 to 2025.

We evaluate model predictions using a rolling-window prediction scheme.
After each match-day $g$, the model is fitted using all matches from rounds 1 to $g$, and predictions are generated for all future matches.
For match-level metrics, we evaluate the predicted probabilities for round $g+1$; for season-level metrics, the model simulates the remaining $380 - g$ matches to produce predictive distributions of accumulated points for each team at every future match-day.
This process begins at round 50 (approximately 5 complete match-days) to ensure sufficient training data, and continues through round 370, yielding predictions for approximately 330 matches per season evaluated out-of-sample.
All reported metrics are averaged across evaluation points within each season and then across the five seasons.

Predictions are generated via Monte Carlo integration over the posterior distribution.
For each future match, 1,000 parameter vectors are sampled from the posterior.
For Bradley-Terry models, match outcomes are sampled directly from the categorical distribution implied by the model.
For Poisson models, home and away goals are simulated independently from Poisson distributions with rates determined by the sampled parameters, and match outcome probabilities are computed as the empirical frequencies of home wins, draws, and away wins.
Point totals are accumulated across simulated matches, yielding predictive distributions from which quantile-based intervals are extracted.

To contextualize model performance, we compare against two baseline predictors:
\begin{itemize}
    \item \textbf{Naive 1}: Assigns uniform probabilities $(1/3, 1/3, 1/3)$ to all outcomes, representing a completely uninformative forecast.
    \item \textbf{Naive 2}: Assigns probabilities based on historical outcome frequencies observed in the current season, reflecting the marginal distribution of results up to match-day $g$.
\end{itemize}

We assess predictive quality using four complementary scoring rules, all negatively oriented (lower values indicate better predictions).
The first three metrics evaluate match-level outcome predictions, while the fourth assesses season-level point forecasts.
Let $n$ denote the number of evaluated matches.
For each match $t$, let $\mathbf{z}_t = (z_{H,t}, z_{D,t}, z_{A,t})$ be the one-hot encoded outcome vector, where $H$, $D$, and $A$ denote home victory, draw, and away victory, respectively.
Let $\mathbf{p}_t = (p_{H,t}, p_{D,t}, p_{A,t})$ be the corresponding predicted probabilities.

The \textbf{Brier Score} (BS) measures the accuracy of probabilistic predictions for categorical outcomes, computed as the mean squared error between predicted probabilities and observed outcomes:
\begin{equation}
    \text{BS} = \frac{1}{n} \sum_{t=1}^{n} \sum_{j \in \{H, D, A\}} (z_{j,t} - p_{j,t})^2.
    \label{eq:brier_score}
\end{equation}

The \textbf{Log Score} (LS), also known as the negative log-likelihood, is a strictly proper scoring rule that heavily penalizes confident but incorrect predictions:
\begin{equation}
    \text{LS} = -\dfrac{1}{n} \sum_{t=1}^{n} \sum_{j \in \{H, D, A\}} z_{j,t} \log(p_{j,t}).
    \label{eq:log_score}
\end{equation}

The \textbf{Ranked Probability Score} (RPS) accounts for the ordinal nature of match outcomes (home win $\succ$ draw $\succ$ away win), penalizing predictions that are ``further'' from the true outcome:
\begin{equation}
    \text{RPS} = \dfrac{1}{2n} \sum_{t=1}^{n} \sum_{k=1}^{2} \left( \sum_{j=1}^{k} (z_{j,t} - p_{j,t}) \right)^2.
    \label{eq:rps}
\end{equation}

Finally, the \textbf{Interval Score} (IS) evaluates the quality of predictive intervals for cumulative points.
Unlike the previous metrics that assess single-match predictions, IS measures how well the model forecasts the range of possible point totals for each team at future match-days.

A good predictive interval should be both \emph{narrow} (providing precise information) and \emph{calibrated} (containing the true value at the stated confidence level).
The IS captures this trade-off: for a confidence level $(1-\alpha)$, let $[l, u]$ be the predicted interval and $y$ the observed points:
\begin{equation}
    \text{IS}_\alpha = (u - l) + \frac{2}{\alpha}(l - y)\mathbf{1}_{y < l} + \frac{2}{\alpha}(y - u)\mathbf{1}_{y > u}.
    \label{eq:interval_score}
\end{equation}
The first term penalizes wide intervals, while the second and third terms penalize coverage failures, which are cases where the observed value falls outside the interval.
The factor $2/\alpha$ ensures that narrower intervals (smaller $\alpha$) incur larger penalties when they fail to cover the true value.

We aggregate IS across all 20 teams, all future match-days, and 19 symmetric prediction intervals ranging from 5\% to 95\% coverage.

For season-level analysis of specific \textit{Brasileirão} seasons, we apply a complementary approach focused on a single well-performing model.
After each match-day, the model is refitted using all available match results, and 10,000 Monte Carlo simulations are used to forecast the remaining matches.
Each iteration samples parameters from the posterior distribution and simulates all unplayed matches according to the model's likelihood, producing a complete final standings table.

The primary output is a full probability matrix assigning each team a probability for each final position (1--20).
From this matrix, we can derive quantities such as:
\begin{itemize}
    \item Probability of winning the championship;
    \item Probability of qualifying for \textit{Copa Libertadores da América} (\textit{Libertadores}), South America's premier continental club competition;
    \item Probability of qualifying for \textit{Copa Sul-Americana} (\textit{Sul-Americana}), South America's secondary continental club competition;
    \item Probability of relegation.
\end{itemize}

Additionally, posterior distributions of team strength parameters are obtained after each match-day, providing interpretable summaries of relative team quality.
